\section{Implementacia}
V tejto časti sa zameriame na samotnú implementáciu navrhovanej aplikácie, ktorá zahŕňa výber technológií, architektúru systému a detailný popis jednotlivých komponentov.

\subsection{Architektonický prístup: MVVM s Jetpack Compose}
V súlade s princípmi čistého návrhu (\textit{Clean Architecture}) opísanými v \cite{sks_medium_2024},
je v projekte implementovaný vzor MVVM s dôrazom na \textbf{jednosmerný tok dát}
(\textit{Unidirectional Data Flow - UDF}). Tradičný mechanizmus \textit{Data Bindingu}
známy zo starších verzií Androidu je tu nahradený moderným reaktívnym prístupom:

\begin{itemize}[noitemsep]
    \item \textbf{State Hoisting (Vynášanie stavu):} ViewModel uchováva kompletný stav obrazovky v podobe imutabilného objektu (napr. \textit{EventsUiState}). UI komponenty sú takmer bezstavové (\textit{stateless}), čo uľahčuje ich znovupoužiteľnosť.
    \item \textbf{Reaktívne pozorovanie (StateFlow):} Používateľské rozhranie je prepojené s \textit{StateFlow} vo ViewModele. Akákoľvek zmena v dátach vyvolá okamžitú \textbf{rekompozíciu} dotknutých častí rozhrania.
    \item \textbf{Spracovanie udalostí (Events):} Všetky akcie používateľa sú prenášané ako udalosti smerom k ViewModelu. View nikdy neupravuje dáta priamo, čo zaručuje, že \textit{ViewModel} zostáva jediným zdrojom pravdy (\textit{Single Source of Truth}).
\end{itemize}

\subsection{Autentifikácia a správa používateľov}
Aplikácia je navrhnutá podľa architektúry \textbf{MVVM (Model-View-ViewModel)}, ktorá zabezpečuje jasné oddelenie biznis logiky od používateľského rozhrania. Implementácia autentifikácie prebieha na úrovni vrstvy \textit{Model} (služba \texttt{AuthService}), zatiaľ čo validácia a spracovanie stavov sú riešené vo vrstve \textit{ViewModel}.

\subsubsection{Implementácia kľúčových funkcií}
Služba \texttt{AuthService} zabezpečuje nízkoúrovňovú komunikáciu s Firebase. Využíva asynchrónne volania pomocou Kotlin korutín, čo umožňuje neblokujúce vykonávanie sieťových operácií. Registrácia nového používateľa zahŕňa nielen vytvorenie účtu, ale aj okamžitú aktualizáciu profilu o zobrazované meno (\textit{displayName}):

\medskip % Pridaná medzera pred kódom
\begin{lstlisting}[language=Kotlin, caption=Registrácia používateľa v AuthService]
suspend fun registerWithEmail(
    email: String,
    password: String,
    displayName: String
): FirebaseUser? {
    val credential = auth.createUserWithEmailAndPassword(email, password).await()
    credential.user?.updateProfile(
        userProfileChangeRequest { this.displayName = displayName }
    )?.await()
    return credential.user
}
\end{lstlisting}
\medskip % Pridaná medzera za kódom

Prihlásenie cez Google účet prebieha prostredníctvom \texttt{GoogleSignInClient}, pričom po získaní \texttt{idToken} sa volá metóda \texttt{firebaseAuthWithGoogle}:

\medskip
\begin{lstlisting}[language=Kotlin, caption=Prihlásenie cez Google v AuthService]
suspend fun firebaseAuthWithGoogle(idToken: String): FirebaseUser? {
    val credential = GoogleAuthProvider.getCredential(idToken, null)
    val result = auth.signInWithCredential(credential).await()
    return result.user
}
\end{lstlisting}
\medskip

Okrem externých poskytovateľov je implementované aj štandardné prihlásenie pomocou e-mailu a hesla:

\medskip
\begin{lstlisting}[language=Kotlin, caption=Prihlásenie pomocou e-mailu a hesla]
suspend fun signInWithEmail(email: String, password: String) {
    auth.signInWithEmailAndPassword(email, password).await()
}
\end{lstlisting}
\medskip

\subsubsection{Validácia vstupných údajov}
Pred odoslaním požiadavky na server vykonáva \texttt{ViewModel} lokálnu validáciu na strane klienta. Týmto prístupom šetríme sieťovú prevádzku a zlepšujeme odozvu rozhrania. V rámci bezpečnosti sa kontroluje:

\begin{itemize}[noitemsep]
    \item formát e-mailovej adresy (validita syntaxe),
    \item komplexnosť hesla (minimálna dĺžka 6 znakov, prítomnosť veľkého písmena a špeciálneho znaku),
    \item zhoda hesla s jeho potvrdením pri registrácii.
\end{itemize}

\medskip
\begin{lstlisting}[language=Kotlin, caption=Validácia hesla v RegisterViewModel]
passwordError = when {
    password.length < 6 -> "Heslo musí mať aspoň 6 znakov."
    !password.any { it.isUpperCase() } -> "Heslo musí obsahovať veľké písmeno."
    !password.any { !it.isLetterOrDigit() } -> "Heslo musí obsahovať špeciálny znak."
    else -> null
}
\end{lstlisting}
\medskip

\subsubsection{Bezpečnostné obmedzenia a banovanie}
Nad rámec štandardnej autentifikácie aplikácia implementuje kontrolu stavu používateľa. Pri každom prihlásení sa overuje atribút \texttt{banned} v dokumente používateľa v databáze Firestore. V prípade, že je používateľ zablokovaný, systém ho okamžite odhlási a zamedzí prístup k chránenému obsahu aplikácie.
\subsection{Implementácia mapového modulu a geolokalizácie}

Ako základ pre mapu je použitá knižnica \textit{Google Maps SDK}, ktorá poskytuje všetky potrebné funkcionality pre prácu s geografickými dátami. Používateľ na hlavnej obrazovke vidí interaktívnu mapu, na ktorej sú jednotlivé udalosti zobrazené vo forme markerov.

\subsubsection{Zobrazenie mapy a integrácia s Compose}
Keďže \texttt{MapView} nie je natívny komponent prostredia \textit{Jetpack Compose}, v aplikácii je využitý \texttt{AndroidView}, ktorý slúži ako most (\textit{bridge}) medzi klasickým \textit{Android View} systémom a \textit{Compose UI}. Mapa sa pri prvotnom načítaní predvolene centruje na súradnice Bratislavy.

\begin{lstlisting}[language=Kotlin, caption=Inicializácia MapView v prostredí Compose]
val mapView = remember { MapView(context).also { it.onCreate(null) } }

AndroidView(factory = { mapView }, modifier = Modifier.fillMaxSize()) { mv ->
    mv.getMapAsync { map ->
        map.moveCamera(CameraUpdateFactory.newLatLngZoom(BRATISLAVA, 13f))
    }
}
\end{lstlisting}

\texttt{MapView} disponuje vlastným životným cyklom, ktorý je potrebné manuálne prepojiť s životným cyklom obrazovky pomocou \texttt{DisposableEffect}. Tým sa zabezpečí, že mapa správne reaguje na udalosti ako \textit{onPause}, \textit{onResume} a \textit{onDestroy}.

\subsubsection{GPS poloha a správa oprávnení}
Aplikácia vyžaduje oprávnenie \texttt{ACCESS\_FINE\_LOCATION}, po ktorého udelení získa poslednú známu polohu zariadenia pomocou \texttt{FusedLocationProviderClient}. Kamera mapy sa presunie na polohu používateľa iba raz pri inicializácii. Tento stav je riadený príznakom \texttt{mapMovedToUser} v \texttt{AppViewModel}, čím sa zabraňuje nežiaducemu a opakovanému presunu kamery pri každej rekomponzícii rozhrania.

\begin{lstlisting}[language=Kotlin, caption=Získanie aktuálnej polohy používateľa]
LocationServices.getFusedLocationProviderClient(context)
    .lastLocation
    .addOnSuccessListener { loc ->
        userLocation = LatLng(loc.latitude, loc.longitude)
        appViewModel.updateUserLocation(loc.latitude, loc.longitude)
    }
\end{lstlisting}

\subsubsection{Práca s markermi a interakcia}
Všetky udalosti sú konvertované na objekty \texttt{MarkerOptions} pomocou metódy \texttt{toMarkerOptions()}, ktorá je definovaná priamo v dátovom modeli \texttt{Event}.

\begin{lstlisting}[language=Kotlin, caption=Konverzia modelu Event na MarkerOptions]
fun toMarkerOptions(icon: BitmapDescriptor? = null): MarkerOptions =
    MarkerOptions()
        .position(LatLng(latitude, longitude))
        .title(title)
        .apply { icon?.let { icon(it) } }
\end{lstlisting}

Zobrazované markery sú synchronizované so stavom \texttt{appViewModel.activeEvents} prostredníctvom \texttt{snapshotFlow}. Pri každej zmene zoznamu udalostí sa mapa reaktívne aktualizuje a pridá nové body s vlastným grafickým vizuálom (čierny pin). Do vlastnosti \texttt{tag} každého markera sa ukladá unikátne \texttt{event.id}, čo umožňuje presnú identifikáciu udalosti pri interakcii. Po kliknutí na marker aplikácia identifikuje udalosť a naviguje používateľa na obrazovku s detailom.

\paragraph{Automatické skrývanie skončených udalostí:}
\texttt{AppViewModel} vystavuje computed property \texttt{activeEvents}, ktorá
filtruje interné pole \texttt{\_events} a vracia len udalosti, ktorých koniec
je v budúcnosti:

\begin{lstlisting}[language=Kotlin, caption=Computed property activeEvents v AppViewModel]
val activeEvents: List<Event>
    get() = _events.values.filter { event ->
        val end = event.dateTo?.time ?: event.dateFrom?.time
        end == null || end > System.currentTimeMillis()
    }
\end{lstlisting}

Pravidlá filtrovania sú:
\begin{itemize}
    \item Ak udalosť má \texttt{dateTo}, rozhoduje tento čas.
    \item Ak \texttt{dateTo} neexistuje, použije sa \texttt{dateFrom}.
    \item Udalosti bez dátumu sú vždy zobrazené.
\end{itemize}

Markery na mape aj filter sheet používajú výhradne \texttt{activeEvents} —
skončené udalosti zmiznú z UI automaticky bez potreby manuálneho mazania
z databázy. Interné \texttt{\_events} zostáva kompletné, aby \texttt{checkNoShows()}
stále videl aj minulé eventy pri kontrole neúčastí.

\textit{Dizajnové rozhodnutie:} Udalosti sa fyzicky nemažú z Firestore — zostávajú
pre históriu účasti, no-show kontrolu a prípadné štatistiky. Skrývanie prebieha
čisto na úrovni UI.

\subsubsection{Okruh pôsobnosti a filtrovanie}
Pre lepšiu prehľadnosť aplikácia filtruje udalosti na základe ich vzdialenosti od používateľa. Markery nachádzajúce sa mimo nastaveného okruhu sú automaticky skryté (\texttt{marker.isVisible = false}). Tento okruh je na mape vizualizovaný ako žltý polopriesvitný kruh (\textit{Circle}), ktorého polomer je prepojený s nastaveniami v \texttt{AppViewModel}.

\subsubsection{Personalizácia zobrazenia}
Používateľ má možnosť prepínať medzi piatimi štýlmi mapového podkladu: \textit{NORMAL, SATELLITE, HYBRID, TERRAIN} a špeciálnym \textit{DARK} režimom, ktorý využíva vlastnú definíciu štýlu vo formáte JSON.
\subsection{Eventy — implementácia}

\subsubsection{Tvorba eventu}

Obrazovka tvorby eventu (\texttt{EventCreationScreen.kt}) sa otvorí po potvrdení
polohy na mape. Formulár obsahuje nasledujúce polia:

\begin{itemize}
    \item \textbf{Názov a popis} — textové polia, obe sú povinné.
    \item \textbf{Miesto} — predvyplnené adresou získanou z Geocoderu / Places API.
    \item \textbf{Cena a počet účastníkov} — hodnota 0 znamená zadarmo, resp. neobmedzený počet.
    \item \textbf{Dátum a čas} — výber cez \texttt{DatePickerDialog} a \texttt{TimePickerDialog}.
    \item \textbf{Trvanie} — zadávané v hodinách a minútach; hodnota \texttt{dateTo}
          sa vypočíta automaticky:
\begin{lstlisting}[language=Kotlin]
val computedDateTo = dateFrom?.let {
    Date(it.time + totalMinutes * 60 * 1000L)
}
\end{lstlisting}
    \item \textbf{Kategória a podkategória} — výber z rozbaľovacieho zoznamu
          definovaného v \texttt{CATEGORY\_MAP}.
    \item \textbf{Viditeľnosť} — verejné alebo súkromné; súkromný event vyžaduje heslo.
    \item \textbf{Typ check-in} — jedna zo štyroch možností:
          \texttt{none}, \texttt{qr}, \texttt{manual}, \texttt{qr\_and\_manual}.
    \item \textbf{Deadline odhlásenia} — počet hodín pred eventom, po uplynutí ktorých
          sa odhlásenie penalizuje.
\end{itemize}

\paragraph{Upload obrázkov cez Cloudinary:}
Obrázky sa vyberajú z galérie zariadenia pomocou
\texttt{ActivityResultContracts.GetMultipleContents()}. Každý obrázok sa uploaduje
ako \texttt{multipart/form-data} priamo na Cloudinary API bez potreby vlastného
backendu:

\begin{lstlisting}[language=Kotlin, caption=Upload obrázka na Cloudinary]
private suspend fun uploadToCloudinary(
    context: Context,
    uri: Uri,
    eventId: String
): String {
    val url = URL("https://api.cloudinary.com/v1_1/$cloudName/image/upload")
    val conn = (url.openConnection() as HttpsURLConnection).apply {
        requestMethod = "POST"
        setRequestProperty(
            "Content-Type",
            "multipart/form-data; boundary=$boundary"
        )
    }
    // zapíše upload_preset + bytes obrázka
    // vráti secure_url z JSON odpovede
}
\end{lstlisting}

Prístupové kľúče (\texttt{CLOUDINARY\_CLOUD\_NAME}, \texttt{CLOUDINARY\_UPLOAD\_PRESET})
sú uložené v súbore \texttt{local.properties} a sprístupnené cez \texttt{BuildConfig}.

\paragraph{Uloženie eventu:}
Po vyplnení formulára sa event uloží volaním \texttt{onSave()}:

\begin{lstlisting}[language=Kotlin, caption=Uloženie eventu]
onSave(event.copy(
    title               = title.trim(),
    dateFrom            = dateFrom,
    dateTo              = computedDateTo,
    visibility          = if (isPrivate) "private" else "public",
    imageUrls           = existingUrls + uploadedUrls,
    checkInType         = checkInType,
    cancelDeadlineHours = cancelDeadlineInput.toIntOrNull() ?: 0
))
\end{lstlisting}

Metóda \texttt{AppViewModel.saveEventToMap()} uloží event súčasne do in-memory
mapy \texttt{\_events} aj do databázy Firestore:

\begin{lstlisting}[language=Kotlin, caption=Uloženie eventu do Firestore]
fun saveEventToMap(event: Event) {
    _events[event.id] = event
    viewModelScope.launch {
        db.collection("events")
          .document(event.id)
          .set(event.toMap())
          .await()
    }
}
\end{lstlisting}

% -------------------------------------------------------
\subsubsection{Detail eventu}

Obrazovka detailu eventu (\texttt{EventDetailScreen.kt}) je rozdelená do troch tabov.

\paragraph{Tab 0 — Detail:}
Zobrazuje všetky informácie o evente. Kľúčové stavy používateľa sú určené
nasledovne:

\begin{lstlisting}[language=Kotlin, caption=Stavové premenné v detaile eventu]
val isCreator    = event.creatorId == currentUserId
val isAttending  = currentUserId in event.attendees
val isOnWaitlist = currentUserId in event.waitlist
val isFull       = event.participants > 0 &&
                   event.attendees.size >= event.participants
val isBanned     = currentUserId in event.bannedUsers
\end{lstlisting}

Tlačidlá akcie sa dynamicky menia podľa stavu — tvorca eventu vidí možnosti
úpravy a mazania, prihlásený účastník vidí odhlásenie, ostatní používatelia
vidia možnosť prihlásenia alebo zaradenia do poradovníka.

\paragraph{Tab 1 — Účastníci:}
Zobrazuje zoznam účastníkov s ich avatarmi. Tvorca eventu má navyše možnosť:
\begin{itemize}
    \item manuálne označiť príchod účastníka (ak \texttt{checkInType == "manual"}),
    \item odstrániť účastníka z eventu — odstránený účastník je zároveň zaradený
          do zoznamu \texttt{bannedUsers} a prvý používateľ z poradovníka je
          automaticky povýšený.
\end{itemize}

\paragraph{Tab 2 — Chat:}
Real-time chat implementovaný pomocou Firestore \texttt{snapshotListener}.
Prístup do chatu majú výlučne prihlásení účastníci a tvorca eventu.

% -------------------------------------------------------
\subsubsection{Správa účasti}

Správa účasti je implementovaná v \texttt{AppViewModel} prostredníctvom
troch metód.

\paragraph{joinEvent — prihlásenie na event:}
Ak je event plný, používateľ je zaradený do poradovníka; v opačnom prípade
priamo medzi účastníkov:

\begin{lstlisting}[language=Kotlin, caption=Prihlásenie na event]
fun joinEvent(eventId: String, userId: String) {
    val isFull = event.participants > 0 &&
                 event.attendees.size >= event.participants
    val updated = if (isFull)
        event.copy(waitlist  = event.waitlist  + userId)
    else
        event.copy(attendees = event.attendees + userId)
    // uloží do Firestore
}
\end{lstlisting}

\paragraph{leaveEvent — odhlásenie z eventu:}
Pri odhlásení sa automaticky povýši prvý používateľ z poradovníka. Ak sa
používateľ odhlási po uplynutí deadlinu, je mu zaznamenaný no-show:

\begin{lstlisting}[language=Kotlin, caption=Odhlásenie z eventu]
fun leaveEvent(eventId: String, userId: String) {
    val applyPenalty = event.isPastCancelDeadline()
    var updated = event.copy(attendees = event.attendees - userId)
    if (updated.waitlist.isNotEmpty()) {
        val next = updated.waitlist.first()
        updated = updated.copy(
            attendees = updated.attendees + next,
            waitlist  = updated.waitlist.drop(1)
        )
    }
    if (applyPenalty) recordNoShow(userId, eventId)
}
\end{lstlisting}

\paragraph{removeAttendee — odstránenie účastníka tvorcom:}
Odstránený účastník je zároveň zablokovaný pre daný event a prvý z poradovníka
je automaticky povýšený:

\begin{lstlisting}[language=Kotlin, caption=Odstránenie účastníka tvorcom eventu]
fun removeAttendee(eventId: String, attendeeId: String) {
    val updated = event.copy(
        attendees   = event.attendees   - attendeeId,
        bannedUsers = event.bannedUsers + attendeeId
    )
    // + automatické povýšenie z poradovníka
}
\end{lstlisting}

% -------------------------------------------------------
\subsubsection{Filtrovanie udalostí}

Filtrovanie je implementované ako samostatná servisná trieda
\texttt{EventFilterService}, ktorá prijíma kritériá a aplikuje ich na zoznam
eventov.

\paragraph{FilterCriteria:}
Dátová trieda zapuzdruje všetky dostupné filtre:

\begin{lstlisting}[language=Kotlin, caption=Dátová trieda FilterCriteria]
data class FilterCriteria(
    val name           : String? = null,
    val category       : String? = null,
    val subcategory    : String? = null,
    val participants   : Int?    = null,
    val maxPrice       : Double? = null,
    val dateFrom       : Date?   = null,
    val dateTo         : Date?   = null,
    val maxDistanceKm  : Double? = null,
    val refLat         : Double  = 0.0,
    val refLng         : Double  = 0.0
)
\end{lstlisting}

\paragraph{EventFilterService:}
Trieda aplikuje kritériá postupne — ak event nevyhovuje čo i len jednému
z nich, je zo zoznamu vyradený:

\begin{lstlisting}[language=Kotlin, caption=Implementácia EventFilterService]
class EventFilterService(private val criteria: FilterCriteria) {

    fun filter(allEvents: List<Event>): List<Event> =
        allEvents.filter { matches(it) }

    private fun matches(event: Event): Boolean {
        if (!criteria.name.isNullOrEmpty() &&
            !event.title.lowercase().contains(criteria.name.lowercase()))
            return false

        if (!criteria.category.isNullOrEmpty() &&
            event.category != criteria.category)
            return false

        if (!criteria.subcategory.isNullOrEmpty() &&
            event.subcategory != criteria.subcategory)
            return false

        if (criteria.participants != null && criteria.participants > 0 &&
            event.participants < criteria.participants)
            return false

        if (criteria.maxPrice != null && criteria.maxPrice < 500.0 &&
            event.price > criteria.maxPrice)
            return false

        if (criteria.dateFrom != null && event.dateFrom != null &&
            event.dateFrom.before(criteria.dateFrom))
            return false

        if (criteria.maxDistanceKm != null) {
            val dist = haversineKm(
                criteria.refLat, criteria.refLng,
                event.latitude,  event.longitude
            )
            if (dist > criteria.maxDistanceKm) return false
        }
        return true
    }
}
\end{lstlisting}

\paragraph{Výpočet vzdialenosti — Haversine:}
Na výpočet vzdialenosti medzi dvoma GPS súradnicami sa používa Haversine vzorec,
ktorý zohľadňuje zakrivenie Zeme. Výsledkom je vzdialenosť v kilometroch
(polomer Zeme $R = 6371\ \text{km}$):

\begin{lstlisting}[language=Kotlin, caption=Haversine vzorec]
private fun haversineKm(
    lat1: Double, lon1: Double,
    lat2: Double, lon2: Double
): Double {
    val dLat = Math.toRadians(lat2 - lat1)
    val dLon = Math.toRadians(lon2 - lon1)
    val a = sin(dLat / 2).pow(2) +
            cos(Math.toRadians(lat1)) *
            cos(Math.toRadians(lat2)) *
            sin(dLon / 2).pow(2)
    return 6371.0 * 2 * asin(sqrt(a))
}
\end{lstlisting}

\paragraph{Filter bottom sheet — UI:}
Filter sa zobrazí ako \texttt{ModalBottomSheet} po kliknutí na tlačidlo
„Nájsť udalosť". Používateľ môže filtrovať podľa:
\begin{itemize}
    \item názvu — fulltextové vyhľadávanie (case-insensitive),
    \item kategórie a podkategórie — výber z \texttt{CATEGORY\_MAP},
    \item počtu osôb — slider v rozsahu 0–200,
    \item maximálnej ceny — slider 0–150\,€; hodnota nad 150\,€ znamená
          bez obmedzenia,
    \item vzdialenosti — slider 1–50\,km od GPS polohy používateľa;
          hodnota nad 50\,km znamená kdekoľvek,
    \item dátumu začiatku — výber cez \texttt{DatePickerDialog}
          a \texttt{TimePickerDialog}.
\end{itemize}

Po kliknutí na „Použiť" sa zostavia kritériá a filtrujú sa markery na mape:

\begin{lstlisting}[language=Kotlin, caption=Aplikovanie filtrov na mapu]
val criteria = FilterCriteria(
    name          = name.trim().ifEmpty { null },
    category      = selectedCategory.ifEmpty { null },
    maxPrice      = if (maxPrice >= 150f) Double.MAX_VALUE
                    else maxPrice.toDouble(),
    maxDistanceKm = if (maxDistance < 50f) maxDistance.toDouble()
                    else null,
    refLat        = userLocation?.latitude  ?: 0.0,
    refLng        = userLocation?.longitude ?: 0.0
)
val filtered = EventFilterService(criteria).filter(events)
markers.values.forEach { it.isVisible = false }
filtered.forEach { event -> markers[event.id]?.isVisible = true }
\end{lstlisting}

Výsledok filtra sa prejaví priamo na mape — zobrazia sa len markery eventov
spĺňajúcich všetky zadané kritériá. Stav filtra sa ukladá do
\texttt{AppViewModel.filterState}, čím sa zachováva pri navigácii medzi
obrazovkami.
\subsection{Hodnotenie udalostí}                                                                                                                                                                                                 
                                                                                                                                                                                                                                   
  Systém hodnotenia umožňuje účastníkom ohodnotiť udalosť po jej skončení.                                                                                                                                                         
  Hodnotenie je implementované priamo v \texttt{AppViewModel}.                                                                                                                                                                                                                                                                                                                                                                                                        
  \subsubsection{Dátový model:}                                                                                                                                                                                                    
  Hodnotenie je uložené priamo v dokumente eventu v troch poliach:                                                                                                                                                                                                                                                                                                                                                                                                    
  \begin{lstlisting}[language=Kotlin, caption=Polia hodnotenia v dátovom modeli eventu]                                                                                                                                            
  val rating      : Double       = 0.0         // vážený priemer
  val totalRatings: Int          = 0           // počet hodnotení
  val ratedBy     : List<String> = emptyList() // UIDs ktorí už hodnotili
  \end{lstlisting}

  \subsubsection{Logika hodnotenia:}
  Nový priemer sa vypočíta ako vážený priemer — starý priemer sa vynásobí počtom
  predchádzajúcich hodnotení, pripočíta sa nové hodnotenie a výsledok sa vydelí
  novým celkovým počtom. Týmto spôsobom sa zachová presnosť bez nutnosti ukladať



  \subsubsection{Pravidlá hodnotenia:}
  \begin{itemize}[noitemsep, nolistsep]
      \item Hodnotiť môžu len účastníci eventu - nie tvorca ani cudzí používatelia.
      \item Každý používateľ môže hodnotiť len raz - kontrola prebieha cez zoznam
            \texttt{ratedBy}.
      \item Hodnotenie je v rozsahu 1 až 5 hviezd.
      \item Priemerné hodnotenie vstupuje ako jeden z faktorov do odporúčacieho
            systému s váhou 16 ak ju používateľ nezmení\,\%.
  \end{itemize}

  \subsubsection{Používateľské rozhranie:}
  Sekcia hodnotenia sa v \texttt{EventDetailScreen} zobrazí len
  účastníkovi, ktorý ešte nehodnotil. Aktuálne priemerné hodnotenie a celkový
  počet hodnotení sú viditeľné pre všetkých používateľov.

  \iffalse                                                                                                                                                                                                                         
  \begin{lstlisting}[language=Kotlin, caption=Podmienka zobrazenia sekcie hodnotenia]
  val isAttending = currentUserId in event.attendees
  val hasRated    = currentUserId in event.ratedBy

  if (isAttending && !isCreator && !hasRated) {
  // 5 hviezd - Row s ikonami
  }
  \end{lstlisting}
  \fi                                                                                                                                                                                                                              

  % -------------------------------------------------------                                                                                                                                                                        
  \subsubsection{Check-in systém}

  Check-in systém umožňuje organizátorovi overiť fyzickú prítomnosť účastníkov
  na udalosti. Pri tvorbe eventu si organizátor zvolí jeden zo štyroch typov
  check-inu, ktoré sú zhrnuté v tabuľke~\ref{tab:checkin-types}.

  \begin{table}[h]
      \centering                                                                                                                                                                                                                   
      \caption{Typy check-inu}
      \label{tab:checkin-types}
      \begin{tabular}{ll}
          \toprule                                                                                                                                                                                                                 
          \textbf{Hodnota} & \textbf{Popis} \\                                                                                                                                                                                     
          \midrule                                                                                                                                                                                                                 
          \texttt{none}            & Check-in sa nevyžaduje \\                                                                                                                                                                     
          \texttt{qr}              & Účastník naskenuje QR kód \\                                                                                                                                                                  
          \texttt{manual}          & Organizátor manuálne označí prítomnosť \\                                                                                                                                                     
          \texttt{qr\_and\_manual} & Kombinácia oboch predchádzajúcich \\                                                                                                                                                          
          \bottomrule                                                                                                                                                                                                              
      \end{tabular}
  \end{table}

  \iffalse                                                                                                                                                                                                                         
  \paragraph{Generovanie QR kódu:}
  QR kód sa generuje pomocou knižnice ZXing priamo na zariadení organizátora
  (\texttt{QrCodeScreen.kt}). Obsahom QR kódu je \texttt{eventId} — unikátny
  identifikátor udalosti z Firestore:

  \begin{lstlisting}[language=Kotlin, caption=Generovanie QR kódu pomocou ZXing]
  private fun generateQrBitmap(content: String, size: Int = 512): Bitmap? {
      val bitMatrix = QRCodeWriter()
          .encode(content, BarcodeFormat.QR_CODE, size, size)
      val bitmap = Bitmap.createBitmap(size, size, Bitmap.Config.RGB_565)
      for (x in 0 until size) {
          for (y in 0 until size) {
              bitmap.setPixel(x, y,
                  if (bitMatrix[x, y]) Color.BLACK else Color.WHITE
              )
          }
      }
      return bitmap
  }
  \end{lstlisting}
  \fi                                                                                                                                                                                                                              

  \paragraph{Generovanie QR kódu:}
  QR kód sa generuje pomocou knižnice ZXing priamo na zariadení organizátora
  (\texttt{QrCodeScreen.kt}). Obsahom kódu je \texttt{eventId} - unikátny
  identifikátor udalosti z Firestore.

  \paragraph{Skenovanie QR kódu:}
  Skenovanie je implementované v \texttt{QrScanScreen.kt} pomocou CameraX
  na zobrazenie náhľadu kamery a ML Kit Barcode Scanning na rozpoznanie kódu.
  Obrazovka najprv požiada o oprávnenie . Stratégia
  \texttt{STRATEGY\_KEEP\_ONLY\_LATEST} zabezpečuje spracovanie len posledného
  snímku. Po úspešnom naskenovaní sa zavolá
  \texttt{appViewModel.checkInAttendee(eventId, userId)}.

  \iffalse                                                                                                                                                                                                                         
  \begin{lstlisting}[language=Kotlin, caption=Skenovanie QR kódu pomocou ML Kit]
  val barcodeScanner = BarcodeScanning.getClient(
      BarcodeScannerOptions.Builder()
          .setBarcodeFormats(Barcode.FORMAT_QR_CODE)
          .build()
  )

  val analysis = ImageAnalysis.Builder()
      .setBackpressureStrategy(ImageAnalysis.STRATEGY_KEEP_ONLY_LATEST)
      .build()

  analysis.setAnalyzer(ContextCompat.getMainExecutor(ctx)) { imageProxy ->
      val image = InputImage.fromMediaImage(
          mediaImage, imageProxy.imageInfo.rotationDegrees
      )
      barcodeScanner.process(image)
          .addOnSuccessListener { barcodes ->
              barcodes.firstOrNull()?.rawValue?.let { value ->
                  scanned = true
                  onScanResult(value)  // value = eventId
              }
          }
  }
  \end{lstlisting}
  \fi                                                                                                                                                                                                                              

  \paragraph{Check-in logika:}
  Metóda \texttt{checkInAttendee()} overí, či je používateľ účastníkom eventu
  a či ešte nebol check-inovaný:

  \begin{lstlisting}[language=Kotlin, caption=Logika check-inu v AppViewModel]
  fun checkInAttendee(eventId: String, userId: String) {
      val event = _events[eventId] ?: return
      if (userId !in event.attendees) {
          checkInError = "Nie si účastníkom tejto udalosti."
          return
      }
      if (userId in event.checkedIn) {
          checkInError = "Už si zaregistrovaný."
          return
      }
      val updated = event.copy(checkedIn = event.checkedIn + userId)
      _events[eventId] = updated
      viewModelScope.launch {
          db.collection("events").document(eventId)
              .update("checkedIn", updated.checkedIn).await()
          checkInSuccess = true
      }
  }
  \end{lstlisting}

  Organizátor môže check-in aj zrušiť volaním \texttt{checkOutAttendee()},
  ktorá odstráni \texttt{userId} zo zoznamu \texttt{checkedIn}.

  \paragraph{Manuálny check-in:}
  Pri type \texttt{manual} alebo \texttt{qr\_and\_manual} vidí organizátor
  v Tab~Účastníci pri každom účastníkovi prepínač $\circ$\,/\,$\checkmark$.
  Kliknutím sa zavolá \texttt{checkInAttendee()} alebo \texttt{checkOutAttendee()}.
  Účastník s potvrdenou prítomnosťou má zelenú odznak „$\checkmark$ Prítomný".

  \paragraph{Väzba na penalizačný systém:}
  Ak je \texttt{checkInType} \texttt{none} a účastník sa neodhlási ani
  nezaregistruje check-in, po skončení eventu mu je automaticky zaznamenaný
  no-show.

  % -------------------------------------------------------                                                                                                                                                                        
  \subsubsection{Systém penalizácií (No-show)}

  Penalizačný systém motivuje používateľov dodržiavať záväzky voči udalostiam.
  Ak sa používateľ prihlási na udalosť s povoleným check-inom, fyzicky sa
  nezúčastní a ani sa včas neodhlási, systém mu zaznamená no-show.

  \paragraph{Dátový model:}
  Penalizácie sú uložené priamo v profile používateľa (\texttt{UserProfile}):

  \begin{lstlisting}[language=Kotlin, caption=Polia no-show a soft ban v UserProfile]
  val noShowTimestamps: List<Long>   = emptyList() // čas každého no-show (epoch ms)
  val noShowEventIds  : List<String> = emptyList() // dedup - eventy kde bol no-show

  val banned          : Boolean      = false        // true = účet zablokovaný
  val banReason       : String       = ""           // "no_show" | "reports" | ""
  val banTimestamp    : Long         = 0L           // čas banu v epoch ms
  \end{lstlisting}

  Systém implementuje \textit{soft ban} --- pri zablokovaní sa používateľské dáta
  (história, priatelia, správy) nemažú. Administrátor môže priamo vo Firebase
  Console vidieť dôvod (\texttt{banReason}) a čas (\texttt{banTimestamp}) banu
  a manuálne rozhodnúť o odbanovaní alebo permanentnom vymazaní účtu.

  \iffalse                                                                                                                                                                                                                         
  Pomocná metóda počíta no-show len za posledných 30 dní:

  \begin{lstlisting}[language=Kotlin, caption=Výpočet počtu no-show za posledných 30 dní]
  fun recentNoShowCount(): Int {
      val cutoff = System.currentTimeMillis() - 30L * 24 * 3_600_000
      return noShowTimestamps.count { it > cutoff }
  }
  \end{lstlisting}
  \fi                                                                                                                                                                                                                              

  \paragraph{Dva spôsoby vzniku no-show:}
  \begin{enumerate}
      \iffalse                                                                                                                                                                                                                     
      \item \textbf{Neskoré odhlásenie} — používateľ sa odhlási po uplynutí
            \texttt{cancelDeadlineHours}:
  \begin{lstlisting}[language=Kotlin]
  fun leaveEvent(eventId: String, userId: String) {
      val applyPenalty = event.isPastCancelDeadline()
      // ... odhlásenie z eventu ...
      if (applyPenalty) recordNoShow(userId, eventId)
  }
  \end{lstlisting}
      \fi                                                                                                                                                                                                                          
      \item \textbf{Neskoré odhlásenie} — používateľ sa odhlási po uplynutí
            \texttt{cancelDeadlineHours} (podrobnejšie popísané v sekcii
            správy účasti).

      \item \textbf{Neúčasť} — systém pri štarte skontroluje všetky minulé
            eventy, kde bol používateľ prihlásený, ale nebol check-inovaný:
  
  \end{enumerate}


  \paragraph{Zablokovanie účtu:}
  Ak ná používateľ 3 zablokovania dostáva ban. Pri ďalšom prihlásení \texttt{loadUserProfile()} toto
  pole skontroluje a okamžite odhlási používateľa — rovnaký mechanizmus je
  popísaný v sekcii autentifikácie.

  \iffalse                                                                                                                                                                                                                         
  \begin{lstlisting}[language=Kotlin, caption=Kontrola zablokovania pri prihlásení]
  if (profile.banned) {
      FirebaseAuth.getInstance().signOut()
      loginBannedError = true
      return@launch
  }
  \end{lstlisting}
  \fi                                                                                                                                                                                                                              

  No-show sa počítajú výlučne za posledných 30 dní — po uplynutí tohto obdobia
  sa staré záznamy nezapočítavajú, čo umožňuje rehabilitáciu používateľa bez
  nutnosti manuálneho zásahu administrátora.

  % -------------------------------------------------------                                                                                                                                                                        
  \subsubsection{Systém nahlasovania}

  Systém nahlasovania umožňuje používateľom upozorniť na nevhodný obsah -
  či už udalosť alebo iného používateľa. Pri dosiahnutí stanoveného počtu
  nahlásení sa obsah automaticky odstráni bez nutnosti zásahu administrátora.

  \iffalse                                                                                                                                                                                                                         
  \paragraph{Dátový model:}
  Každé nahlásenie je reprezentované triedou \texttt{Report} a uložené
  do Firestore kolekcie \texttt{reports}:

  \begin{lstlisting}[language=Kotlin, caption=Dátová trieda Report]
  data class Report(
      val id        : String = "",
      val reporterId: String = "",
      val targetId  : String = "",
      val targetType: String = "",
      val reason    : String = "",
      val timestamp : Long   = System.currentTimeMillis()
  )
  \end{lstlisting}
  \fi                                                                                                                                                                                                                              

  Každé nahlásenie je uložené do Firestore kolekcie \texttt{reports} s poliami
  \texttt{reporterId}, \texttt{targetId}, \texttt{targetType} (\texttt{"event"}
  alebo \texttt{"user"}), \texttt{reason} a \texttt{timestamp}.

  \paragraph{Dôvody nahlasovania:}
  Zoznam dôvodov a hranica automatickej akcie sú definované ako konštanty
  v \texttt{AppViewModel}:

  \begin{lstlisting}[language=Kotlin, caption=Konštanty systému nahlasovania]
  companion object {
      val REPORT_REASONS = listOf(
          "Nevhodný obsah",
          "Spam",
          "Falošná udalosť / profil",
          "Obťažovanie",
          "Iné"
      )
      private const val REPORTS_TO_BAN = 3
  }
  \end{lstlisting}

 
  


  \paragraph{Logika nahlasovania používateľa:}
  Metóda \texttt{reportUser()} funguje rovnakým spôsobom. Pri dosiahnutí
  prahu sa namiesto mazania udalosti nastaví príznak zablokovania
  na profile používateľa:

  \begin{lstlisting}[language=Kotlin, caption=Automatické zablokovanie používateľa]
  if (total.size() >= REPORTS_TO_BAN) {
      val banTime = System.currentTimeMillis()
      db.collection("users")
        .document(targetUserId)
        .update(mapOf(
            "banned"       to true,
            "banReason"    to "reports",
            "banTimestamp" to banTime
        )).await()
      cancelPendingInvitations(targetUserId)
  }
  \end{lstlisting}

  \iffalse                                                                                                                                                                                                                         
  \paragraph{Zrušenie čakajúcich pozvánok pri bane:}
  Pri každom zablokovaní účtu sa automaticky vymažú všetky čakajúce pozvánky
  adresované zabanovanému používateľovi:

  \begin{lstlisting}[language=Kotlin, caption=Implementácia cancelPendingInvitations]
  suspend fun cancelPendingInvitations(userId: String) {
      val pending = db.collection("invitations")
          .whereEqualTo("receiverId", userId)
          .whereEqualTo("status", "pending")
          .get().await()
      pending.documents.forEach { it.reference.delete().await() }
  }
  \end{lstlisting}
  \fi                                                                                                                                                                                                                              

  Pri každom zablokovaní (no-show aj nahlásenia) sa automaticky vymažú všetky
  čakajúce pozvánky adresované zabanovanému používateľovi volaním
  \texttt{cancelPendingInvitations()}.

  \paragraph{Pravidlá systému:}
  Tabuľka~\ref{tab:report-rules} sumarizuje pravidlá platné pre systém
  nahlasovania.

  \begin{table}[h]
      \centering                                                                                                                                                                                                                   
      \caption{Pravidlá systému nahlasovania}
      \label{tab:report-rules}
      \begin{tabular}{lp{8.5cm}}
          \toprule                                                                                                                                                                                                                 
          \textbf{Pravidlo} & \textbf{Detail} \\                                                                                                                                                                                   
          \midrule                                                                                                                                                                                                                 
          Kto môže nahlásiť &                                                                                                                                                                                                      
              Každý prihlásený používateľ okrem tvorcu udalosti
              a seba samého. \\                                                                                                                                                                                                    
          Duplikát &                                                                                                                                                                                                               
              Jeden používateľ môže nahlásiť ten istý cieľ len raz. \\                                                                                                                                                             
          Prah &                                                                                                                                                                                                                   
              3 unikátne nahlásenia spustia automatickú akciu. \\                                                                                                                                                                  
          Udalosť &                                                                                                                                                                                                                
              Vymazaná z Firestore aj z in-memory mapy. \\                                                                                                                                                                         
          Používateľ &                                                                                                                                                                                                             
              \texttt{banned = true} $\rightarrow$ odhlásenie
              pri ďalšom prihlásení. \\                                                                                                                                                                                            
          \bottomrule                                                                                                                                                                                                              
      \end{tabular}
  \end{table}

  % -------------------------------------------------------                                                                                                                                                                        
  \subsection{Implementácia komunikačného modulu v reálnom čase}

  Základom interaktivity je integračná vrstva služby Firebase Cloud
  Firestore, ktorá umožňuje obojsmernú synchronizáciu dát v reálnom čase
  bez nutnosti vlastnej serverovej infraštruktúry a správy WebSocket spojení.

  \iffalse                                                                                                                                                                                                                         
  \subsubsection{Dátový model správy}
  Pre potreby chatu bol definovaný dátový model \texttt{ChatMessage}:

  \begin{lstlisting}[language=Kotlin, caption={Dátový model správy v jazyku Kotlin}]
  data class ChatMessage(
      val id: String = "",
      val senderId: String = "",
      val senderName: String = "",
      val text: String = "",
      val timestamp: Long = 0L
  )
  \end{lstlisting}
  \fi                                                                                                                                                                                                                              

  Správy sú uložené v podkolekcii \texttt{events/\{eventId\}/messages/}
  s poliami \texttt{senderId}, \texttt{senderName}, \texttt{text}
  a \texttt{timestamp}.

  \subsubsection{Reaktívna synchronizácia dát}
  Technické jadro komunikácie tvorí metóda \texttt{addSnapshotListener}.
  Na rozdiel od štandardných HTTP požiadaviek vytvára trvalé spojenie
  s podkolekciou správ. Akonáhle dôjde k zápisu novej správy, Firebase
  iniciuje push notifikáciu a aplikácia aktualizuje reaktívny stav
  \texttt{chatMessages} vo ViewModeli.

  \begin{lstlisting}[language=Kotlin, caption={Implementácia real-time listenera vo ViewModeli}]
  fun startChatListener(eventId: String) {
      chatListener?.remove()
      chatListener = db.collection("events").document(eventId)
          .collection("messages")
          .orderBy("timestamp")
          .addSnapshotListener { snapshot, _ ->
              _chatMessages.clear()
              snapshot?.documents?.forEach { doc ->
                  doc.toObject(ChatMessage::class.java)?.let {
                      _chatMessages.add(it)
                  }
              }
          }
  }
  \end{lstlisting}
\subsubsection{Systém priateľstiev}

Systém priateľstiev umožňuje používateľom nadviazať spojenie a komunikovať
priamo cez správy. Implementácia pokrýva posielanie žiadostí, ich
prijímanie a odmietanie, odstraňovanie priateľov a priamu komunikáciu.

\paragraph{Dátový model:}
Každé priateľstvo je reprezentované triedou \texttt{Friendship} a uložené
do Firestore kolekcie \texttt{friendships}:

\begin{lstlisting}[language=Kotlin, caption=Dátová trieda Friendship]
data class Friendship(
    val id         : String = "",
    val requesterId: String = "",        // kto poslal žiadosť
    val receiverId : String = "",        // kto prijíma žiadosť
    val status     : String = "pending", // "pending" | "accepted"
    val timestamp  : Long   = 0L
)
\end{lstlisting}

\paragraph{Stavy priateľstva:}
Pre každú dvojicu používateľov môže existovať práve jeden vzťah.
Metóda \texttt{getFriendshipStatus()} vráti jeden zo štyroch stavov
uvedených v tabuľke~\ref{tab:friendship-states}:

\begin{lstlisting}[language=Kotlin, caption=Určenie stavu priateľstva]
fun getFriendshipStatus(currentUserId: String, otherUserId: String): String {
    val f = _friendships.firstOrNull {
        (it.requesterId == currentUserId && it.receiverId == otherUserId) ||
        (it.requesterId == otherUserId   && it.receiverId == currentUserId)
    } ?: return "none"

    return when {
        f.status == "accepted"         -> "friends"
        f.requesterId == currentUserId -> "pending_sent"
        else                           -> "pending_received"
    }
}
\end{lstlisting}

\begin{table}[h]
    \centering
    \caption{Stavy priateľstva}
    \label{tab:friendship-states}
    \begin{tabular}{ll}
        \toprule
        \textbf{Stav} & \textbf{Význam} \\
        \midrule
        \texttt{none}             & Žiadny vzťah medzi používateľmi. \\
        \texttt{pending\_sent}     & Aktuálny používateľ poslal žiadosť. \\
        \texttt{pending\_received} & Aktuálny používateľ dostal žiadosť. \\
        \texttt{friends}           & Priateľstvo je akceptované. \\
        \bottomrule
    \end{tabular}
\end{table}

Stav sa využíva v \texttt{PublicProfileScreen} na zobrazenie správneho
tlačidla — „Pridať priateľa", „Žiadosť odoslaná", „Prijať\,/\,Odmietnuť"
alebo „Odstrániť priateľa".

\paragraph{Načítanie priateľstiev:}
Pri prihlásení sa načítajú všetky dokumenty, kde je používateľ buď
odosielateľom alebo prijímateľom. Odmietnuté žiadosti sa nenačítavajú —
zabraňuje sa tým opakovanému posielaniu žiadostí od toho istého
používateľa:

\begin{lstlisting}[language=Kotlin, caption=Načítanie priateľstiev z Firestore]
fun loadFriendships(uid: String) {
    val sent     = db.collection("friendships")
        .whereEqualTo("requesterId", uid).get().await()
    val received = db.collection("friendships")
        .whereEqualTo("receiverId",  uid).get().await()
    (sent.documents + received.documents).forEach { doc ->
        val f = Friendship.fromFirestore(doc.id, doc.data!!)
        if (f.status != "declined") _friendships.add(f)
    }
}
\end{lstlisting}

\paragraph{Odoslanie žiadosti:}
Pred vytvorením novej žiadosti sa skontroluje, či medzi dvojicou
používateľov už vzťah existuje:

\begin{lstlisting}[language=Kotlin, caption=Odoslanie žiadosti o priateľstvo]
fun sendFriendRequest(fromId: String, toId: String) {
    val alreadyExists = _friendships.any {
        (it.requesterId == fromId && it.receiverId == toId) ||
        (it.requesterId == toId   && it.receiverId == fromId)
    }
    if (alreadyExists) {
        friendshipError = "Žiadosť o priateľstvo už existuje."
        return
    }
    val friendship = Friendship(
        id          = docRef.id,
        requesterId = fromId,
        receiverId  = toId,
        status      = "pending",
        timestamp   = System.currentTimeMillis()
    )
    docRef.set(friendship.toMap()).await()
    _friendships.add(friendship)
}
\end{lstlisting}

\paragraph{Odpoveď na žiadosť:}
Pri prijatí sa stav aktualizuje na \texttt{accepted}; pri odmietnutí
sa dokument z Firestore aj z lokálnej pamäte vymaže:

\begin{lstlisting}[language=Kotlin, caption=Prijatie alebo odmietnutie žiadosti]
fun respondToFriendRequest(friendshipId: String, accept: Boolean) {
    if (accept) {
        db.collection("friendships").document(friendshipId)
            .update("status", "accepted").await()
        _friendships[idx] = _friendships[idx].copy(status = "accepted")
    } else {
        db.collection("friendships").document(friendshipId)
            .delete().await()
        _friendships.removeAll { it.id == friendshipId }
    }
}
\end{lstlisting}

\paragraph{Odstránenie priateľa:}
Metóda \texttt{removeFriend()} vyhľadá akceptované priateľstvo medzi
dvojicou a vymaže ho z Firestore aj z lokálnej pamäte:

\begin{lstlisting}[language=Kotlin, caption=Odstránenie priateľa]
fun removeFriend(currentUserId: String, otherUserId: String) {
    val f = _friendships.firstOrNull {
        it.status == "accepted" &&
        ((it.requesterId == currentUserId && it.receiverId == otherUserId) ||
         (it.requesterId == otherUserId   && it.receiverId == currentUserId))
    } ?: return
    _friendships.remove(f)
    db.collection("friendships").document(f.id).delete().await()
}
\end{lstlisting}

\paragraph{Priama komunikácia:}
Priatelia si môžu písať správy cez \texttt{FriendChatScreen}. ID
konverzácie sa generuje deterministicky zo zoradených UID oboch
používateľov — zaručuje sa tak, že obaja vždy pristupujú k tomu istému
dokumentu bez ohľadu na to, kto rozhovor začal:

\begin{lstlisting}[language=Kotlin, caption=Deterministické generovanie ID konverzácie]
private fun directConversationId(uid1: String, uid2: String) =
    if (uid1 < uid2) "${uid1}_${uid2}" else "${uid2}_${uid1}"
\end{lstlisting}

Správy sa ukladajú do Firestore kolekcie
\texttt{directMessages/\{conversationId\}/messages} a synchronizujú sa
v reálnom čase cez \texttt{snapshotListener} — rovnaký mechanizmus
ako chat v udalostiach.

\paragraph{Používateľské rozhranie:}
Obrazovka \texttt{FriendsScreen} je dostupná z bočného menu pod položkou
„Priatelia" a je rozdelená do dvoch sekcií:
\begin{itemize}
    \item \textbf{Žiadosti o priateľstvo} — zobrazujú sa len ak existujú,
          s tlačidlami „Prijať" a „Odmietnuť".
    \item \textbf{Priatelia} — zoznam akceptovaných priateľov s možnosťou
          otvoriť priamu konverzáciu alebo priateľa odstrániť.
\end{itemize}
\subsection{Odporúčací systém}                                                                                                                                                                                                   
                                                                                          
  Pre aplikáciu Joinly sme zvolili Weighted Hybrid System -
  váhový hybridný systém kombinujúci tri prístupy:
                                          
  \begin{itemize}[noitemsep, nolistsep]                                                                                                                                                                                                                 
      \item \textbf{Content-Based} — odporúčania na základe kategórií                                                                                                                                                                          navštívených eventov,                                                                                                                                                                                                  
      \item \textbf{Kontextový} — zohľadňuje vzdialenosť, čas a popularitu,                                                                                                                                                        
      \item \textbf{Explicitná spätná väzba} — bonus za eventy z obľúbených
            kategórií.
  \end{itemize}

  Každý z týchto prístupov prispieva do výsledného skóre s nastaviteľnou
  váhou, ktorú si používateľ môže upraviť priamo v aplikácii pomocou sliderov.

  \paragraph{Prečo hybridný systém?}
  \begin{enumerate}[noitemsep, nolistsep]
      \item \textbf{Kvalitné metadáta} — každý event má kategóriu,
            podkategóriu, lokáciu a hodnotenie.
      \item \textbf{Menšia používateľská báza} — Collaborative filtering
            vyžaduje veľké množstvo používateľov; hybridný prístup funguje
            aj bez nich.
      \item \textbf{Vysvetliteľnosť} — systém dokáže zdôvodniť odporúčanie
            (napr.\ „odporúčame ti toto, lebo máš rád športové eventy").
      \item \textbf{Flexibilita} — používateľ si sám nastaví čo je pre neho
            dôležitejšie (blízkosť vs.\ kategória vs.\ popularita).
      \item \textbf{Real-time výpočet} — skóre sa počíta za behu,
            nevyžadujú sa predpočítané modely.
  \end{enumerate}

  \paragraph{Dvojúrovňové kategórie:}
  Štandardný Content-Based systém používa jednu úroveň kategórií.
  V našom systéme sú kategórie dvojúrovňové:
  \begin{enumerate}[noitemsep, nolistsep]
      \item \textbf{Hlavná kategória} — široká skupina (napr.\ \textit{Sports}),
      \item \textbf{Podkategória} — špecifická (napr.\ \textit{Cycling sports}).
  \end{enumerate}

  Toto umožňuje presnejšie odporúčania — používateľ, ktorý navštívil
  cyklistický event, dostane bonus pre ďalšie cyklistické eventy, nie len
  všeobecne športové. Previazanie medzi úrovňami zabezpečuje trieda
  \textbf{CategoryMapper}.

  % ═══════════════════════════════════════════════════════════════════════                                                                                                                                                        
  \subsubsection{Váhy a výsledné skóre}
  % ═══════════════════════════════════════════════════════════════════════                                                                                                                                                        

  Systém používa 6 faktorov, každý s vlastnou váhou. Predvolené hodnoty
  sú definované v dátovej triede \texttt{RecommendationWeights}:

  \begin{lstlisting}[language=Kotlin, caption=Predvolené váhy odporúčacieho systému]
  data class RecommendationWeights(
      val mainCategory  : Float = 0.20f,  // 20 %                                                                                                                                                                                  
      val subCategory   : Float = 0.16f,  // 16 %                                                                                                                                                                                  
      val distance      : Float = 0.16f,  // 16 %                                                                                                                                                                                  
      val rating        : Float = 0.16f,  // 16 %                                                                                                                                                                                  
      val popularity    : Float = 0.16f,  // 16 %                                                                                                                                                                                  
      val favoriteBonus : Float = 0.16f   // 16 %                                                                                                                                                                                  
  )
  \end{lstlisting}

  Hlavná kategória má mierne vyššiu váhu (20\,\%) pretože je
  najspoľahlivejším indikátorom záujmu používateľa.
   Keďže súčet váh musí byť vždy rovný 1 (100\,\%), systém hodnoty
  automaticky normalizuje po každej zmene slidera.

  \begin{equation}
  \sum_{i=1}^{6} w_i = 1{,}0 \quad \text{(vždy 100\,\%)}
  \end{equation}

  \paragraph{Vzorec celkového skóre:}

  \begin{equation}
  \boxed{
  Score = w_1 \cdot C_{main} + w_2 \cdot C_{sub} + w_3 \cdot D +
          w_4 \cdot R + w_5 \cdot P + w_6 \cdot B_{fav}
  }
  \end{equation}

  Kde:
  \begin{itemize}[noitemsep, nolistsep]
      \item $w_1 \ldots w_6$ — nastaviteľné váhy (súčet = 1,0),
      \item $C_{main}$ — skóre hlavnej kategórie $(0{,}0 - 1{,}0)$,
      \item $C_{sub}$ — skóre podkategórie $(0{,}0 - 1{,}0)$,
      \item $D$ — skóre vzdialenosti $(0{,}0 - 1{,}0)$,
      \item $R$ — skóre ratingu $(0{,}0 - 1{,}0)$,
      \item $P$ — skóre popularity $(0{,}0 - 1{,}0)$,
      \item $B_{fav}$ — bonus za obľúbenú kategóriu $(0{,}0 - 1{,}0)$.
  \end{itemize}

  Keďže každý faktor je v rozsahu $0{,}0 - 1{,}0$ a súčet váh je vždy
  1,0, výsledné skóre je garantovane v rozsahu:

  \begin{equation}
  0{,}0 \leq Score \leq 1{,}0
  \end{equation}

  % ═══════════════════════════════════════════════════════════════════════                                                                                                                                                        
  \subsubsection{Jednotlivé faktory}
  % ═══════════════════════════════════════════════════════════════════════                                                                                                                                                        

  \paragraph{Hlavná kategória ($C_{main}$):}
  Systém zostaví profil používateľa z navštívených eventov - koľko percent
  navštívených eventov patrilo do ktorej kategórie. Skóre sa vypočíta
  nasledovne:

  \begin{lstlisting}[language=Kotlin, caption=Výpočet skóre hlavnej kategórie]
  var mainMatch = userProfile[mainCategory] ?: 0.0
  if (userProfile.isEmpty()) mainMatch = 0.5   // novy pouzivatel
  else if (mainMatch == 0.0) mainMatch = 0.05  // nikdy nenavstivena

  val mainScore = mainMatch * w.mainCategory
  \end{lstlisting}

  \begin{itemize}[noitemsep, nolistsep]
      \item Ak používateľ navštívil 10 eventov, z toho 4 boli „Sports"
            $\rightarrow$ preferencia Sports = 40\,\% $(C_{main} = 0{,}4)$.
      \item Pre nových používateľov (bez histórie) $\rightarrow$                                                                                                                                                                   
            $C_{main} = 0{,}5$ (neutrálne skóre).
      \item Pre nikdy nenavštívenú kategóriu $\rightarrow$                                                                                                                                                                         
            $C_{main} = 0{,}05$ (aby systém nevrátil prázdny zoznam).
  \end{itemize}

  \paragraph{Podkategória ($C_{sub}$):}
  Funguje rovnako ako hlavná kategória, ale na úrovni podkategórií.
  Skóre podkategórie dopĺňa hlavnú kategóriu a umožňuje rozlišovať
  medzi napríklad cyklistickými a plaveckými eventmi v rámci rovnakej
  skupiny Sports.

  \paragraph{Vzdialenosť ($D$):}
  Vzdialenosť sa počíta Haversinovým vzorcom (popísaným v sekcii
  filtrovania) a normalizuje sa tak, aby bližšie eventy dostávali
  vyššie skóre:

  \begin{lstlisting}[language=Kotlin, caption=Výpočet skóre vzdialenosti]
  val distance  = haversineKm(userLat, userLng,
                               event.latitude, event.longitude)
  val distScore = (1.0 / (1.0 + distance / 10.0)) * w.distance
  \end{lstlisting}

  \paragraph{Rating ($R$):}
  Hodnotenie eventu sa normalizuje do rozsahu $0{,}0 - 1{,}0$. Eventy
  bez hodnotení dostanú neutrálne skóre 0,5, čím sa zabraňuje ich
  nespravodlivému penalizovaniu:

  \paragraph{Popularita ($P$):}
  Popularita sa meria ako pomer aktuálneho počtu účastníkov voči
  maximálnej kapacite eventu. Ak event nemá stanovený limit, používa
  sa strop 100:


  \paragraph{Bonus za obľúbenú kategóriu ($B_{fav}$):}
  Obľúbené eventy sa \textbf{nepočítajú} do profilu preferencií —
  používajú sa \textbf{výhradne} na tento bonus. To umožňuje
  používateľovi explicitne označiť čo sa mu páči bez ovplyvnenia
  celkového profilu navštívených eventov:

  % ═══════════════════════════════════════════════════════════════════════                                                                                                                                                        
  \subsubsection{Výstup a používateľské rozhranie}
  % ═══════════════════════════════════════════════════════════════════════                                                                                                                                                        

  Výsledkom výpočtu je zoznam objektov ScoredEvent zoradený
  zostupne podľa skóre. Odporúčania sa zobrazujú na obrazovke
  RecommendedEventsScreen, kde každý event zobrazuje aj
  percentuálne skóre zhody. Používateľ si môže jednotlivé váhy faktorov
  upraviť pomocou sliderov priamo na tejto obrazovke — po každej zmene
  sa odporúčania prepočítajú v reálnom čase.
 \subsection{Implementácia systému pozvánok}                                                                                                                                                                                      

  Systém pozvánok umožňuje organizátorovi udalosti prizvať konkrétnych
  používateľov priamo cez ich e-mailovú adresu. Pozvaný používateľ dostane
  pozvánku v zozname \textit{Pozvánky}, kde ju môže prijať alebo odmietnuť.
                                                                                                                                                                                                                                     \iffalse
  \paragraph{Dátový model:}                                                                                                                                                                                                        
  Každá pozvánka je reprezentovaná objektom \texttt{Invitation} a uložená                                                                                                                                                          
  ako dokument vo Firestore kolekcii \texttt{invitations}. Model uchováva
  identifikátory odosielateľa aj príjemcu, referenciu na udalosť a aktuálny
  stav pozvánky:

  \begin{lstlisting}[language=Kotlin, caption=Dátový model Invitation]
  data class Invitation(
      val id           : String = "",
      val eventId      : String = "",
      val eventTitle   : String = "",
      val eventPlace   : String = "",
      val fromUserId   : String = "",
      val fromUserName : String = "",
      val toUserId     : String = "",
      val toEmail      : String = "",
      val status       : String = "pending",
      val createdAt    : Date   = Date()
  )
  \end{lstlisting}
  \fi                                                                                                                                                                                                                              

  Každá pozvánka je uložená ako dokument vo Firestore kolekcii
  \texttt{invitations} s poliami \texttt{eventId}, \texttt{fromUserId},
  \texttt{toUserId}, \texttt{toEmail}, \texttt{status}
  (\texttt{"pending"} | \texttt{"accepted"} | \texttt{"declined"})
  a \texttt{createdAt}.

  \paragraph{Odoslanie pozvánky:}
  Organizátor odošle pozvánku priamo z detailu udalosti zadaním
  e-mailovej adresy pozývaného. Metóda \texttt{sendInvitation} najprv
  overí existenciu používateľa v kolekcii \texttt{users} a následne
  skontroluje duplicitu — ak pozvánka pre danú kombináciu
  \texttt{eventId} a \texttt{toUserId} už existuje, odoslanie sa zamietne.

  \paragraph{Zobrazenie pozvánok:}
  Obrazovka \texttt{MyInvitationsScreen} rozdeľuje pozvánky do dvoch
  skupín: \textit{Čakajúce} (stav \texttt{pending}) a \textit{Vybavené}
  (stav \texttt{accepted} alebo \texttt{declined}).

  \paragraph{Reakcia na pozvánku:}
  Pri prijatí sa stav aktualizuje na \texttt{accepted} a automaticky
  sa zavolá \texttt{joinEvent}, čím sa používateľ priradí medzi
  účastníkov udalosti. Pri odmietnutí sa stav nastaví na
  \texttt{declined} bez ďalšej akcie.
\subsection{Implementácia systému lokálnych notifikácií}

Aplikácia Joinly informuje používateľov o nadchádzajúcich udalostiach
prostredníctvom lokálnych notifikácií systému Android. Na rozdiel od push
notifikácií, ktoré vyžadujú serverovú infraštruktúru, lokálne notifikácie
sú naplánované priamo na zariadení a nevyžadujú aktívne internetové
pripojenie v momente doručenia.

\paragraph{Architektúra systému:}
Systém pozostáva z troch spolupracujúcich komponentov:
\begin{itemize}
    \item \textbf{NotificationHelper} — singleton objekt zodpovedný za
          plánovanie a rušenie notifikácií,
    \item \textbf{NotificationReceiver} — prijímač systémových správ,
          ktorý zostavuje a zobrazuje notifikáciu,
    \item \textbf{Notification} — dátový model zabezpečujúci perzistenciu
          naplánovaných notifikácií vo Firestore.
\end{itemize}

\paragraph{Notifikačný kanál a správa oprávnení:}
Od verzie Android 8.0 (Oreo) je povinné zaregistrovať notifikačný kanál
pred prvým použitím. Kanál s identifikátorom \texttt{joinly\_events}
a prioritou \texttt{IMPORTANCE\_HIGH} sa vytvára pri štarte aplikácie
v \texttt{LaunchedEffect} v kompozícii \texttt{App.kt}:

\begin{lstlisting}[language=Kotlin, caption=Vytvorenie notifikačného kanála]
fun createChannel(context: Context) {
    if (Build.VERSION.SDK_INT >= Build.VERSION_CODES.O) {
        val channel = NotificationChannel(
            CHANNEL_ID, CHANNEL_NAME, NotificationManager.IMPORTANCE_HIGH
        ).apply {
            description = "Upozornenia na nadchadzajuce udalosti"
        }
        context.getSystemService(NotificationManager::class.java)
            .createNotificationChannel(channel)
    }
}
\end{lstlisting}

Od verzie Android 13 (Tiramisu) vyžaduje zobrazovanie notifikácií
explicitné oprávnenie \texttt{POST\_NOTIFICATIONS}. Aplikácia ho vyžiada
pri prvom spustení pomocou \texttt{rememberLauncherForActivityResult}.
Oprávnenie \texttt{SCHEDULE\_EXACT\_ALARM} je deklarované v súbore
\texttt{AndroidManifest.xml} a umožňuje presné načasovanie budíkov.

\paragraph{Plánovanie notifikácie:}
Notifikácia sa naplánuje automaticky v momente, keď sa používateľ pripojí
k udalosti ako potvrdený účastník (nie ako čakajúci v poradovníku).
Upozornenie sa odošle hodinu pred začiatkom udalosti. Plánovanie
zabezpečuje \texttt{AlarmManager} prostredníctvom metódy
\texttt{setExactAndAllowWhileIdle}, ktorá garantuje spustenie aj v prípade,
že je zariadenie v úspornom režime (\textit{Doze mode}). Na zariadeniach
s Android 12+ sa pred použitím presného alarmu overí dostupnosť oprávnenia:

\begin{lstlisting}[language=Kotlin, caption=Plánovanie notifikácie pomocou AlarmManager]
fun scheduleEventNotification(context: Context, event: Event) {
    val dateFrom = event.dateFrom ?: return
    val notifyAt = dateFrom.time - NOTIFY_BEFORE_MS   // 1 hodina pred
    if (notifyAt <= System.currentTimeMillis()) return

    val alarmManager = context.getSystemService(AlarmManager::class.java)
    val pending = buildPendingIntent(context, event)

    if (Build.VERSION.SDK_INT >= Build.VERSION_CODES.S
            && !alarmManager.canScheduleExactAlarms()) {
        alarmManager.setAndAllowWhileIdle(
            AlarmManager.RTC_WAKEUP, notifyAt, pending
        )
    } else {
        alarmManager.setExactAndAllowWhileIdle(
            AlarmManager.RTC_WAKEUP, notifyAt, pending
        )
    }
}
\end{lstlisting}

Každý alarm je identifikovaný pomocou \texttt{PendingIntent}, do ktorého
sú vložené metadáta udalosti — \texttt{eventId}, \texttt{eventTitle},
\texttt{eventPlace} a čas začiatku. Ako kľúč pre \texttt{requestCode}
sa používa \texttt{eventId.hashCode()}, čo zaručuje jednoznačnú
identifikáciu každej naplánovanej notifikácie.

\paragraph{Zobrazenie notifikácie:}
Keď \texttt{AlarmManager} spustí alarm, systém doručí broadcast triede
\texttt{NotificationReceiver}, ktorá implementuje \texttt{BroadcastReceiver}.
Prijímač z \texttt{Intent} extrahuje metadáta udalosti, zostaví text správy
a zobrazí notifikáciu cez \texttt{NotificationManagerCompat}:

\begin{lstlisting}[language=Kotlin, caption=Zostavenie a zobrazenie notifikácie]
override fun onReceive(context: Context, intent: Intent) {
    val eventTitle = intent.getStringExtra("eventTitle") ?: "Udalost"
    val eventPlace = intent.getStringExtra("eventPlace") ?: ""
    val eventTime  = intent.getLongExtra("eventTime", 0L)

    val timeStr = SimpleDateFormat("HH:mm", Locale.getDefault())
        .format(Date(eventTime))
    val body = "\"$eventTitle\" zacina o $timeStr - $eventPlace"

    val notification = NotificationCompat.Builder(
            context, NotificationHelper.CHANNEL_ID
        )
        .setSmallIcon(R.drawable.ic_star)
        .setContentTitle("Joinly -- Event o hodinu!")
        .setContentText(body)
        .setStyle(NotificationCompat.BigTextStyle().bigText(body))
        .setPriority(NotificationCompat.PRIORITY_HIGH)
        .setAutoCancel(true)
        .build()

    NotificationManagerCompat.from(context)
        .notify(eventId.hashCode(), notification)
}
\end{lstlisting}

Trieda \texttt{NotificationReceiver} je registrovaná v súbore
\texttt{AndroidManifest.xml} ako prijímač systémových správ.

\paragraph{Zrušenie notifikácie:}
Ak používateľ opustí udalosť alebo ak organizátor udalosť vymaže,
naplánovaný alarm sa zruší volaním \texttt{cancelEventNotification}.
Zrušenie sa vykoná vyhľadaním existujúceho \texttt{PendingIntent}
s príznakmi \texttt{FLAG\_NO\_CREATE} — ak alarm neexistuje, metóda
nevyvolá chybu:

\begin{lstlisting}[language=Kotlin, caption=Zrušenie naplánovanej notifikácie]
fun cancelEventNotification(context: Context, eventId: String) {
    val alarmManager = context.getSystemService(AlarmManager::class.java)
    val intent = Intent(context, NotificationReceiver::class.java)
    val pending = PendingIntent.getBroadcast(
        context,
        eventId.hashCode(),
        intent,
        PendingIntent.FLAG_NO_CREATE or PendingIntent.FLAG_IMMUTABLE
    )
    pending?.let { alarmManager.cancel(it) }
}
\end{lstlisting}

\paragraph{Perzistencia vo Firestore:}
Okrem lokálneho alarmu sa každá naplánovaná notifikácia uloží aj do
Firestore kolekcie \texttt{notifications}. Dátový model \texttt{Notification}
uchováva \texttt{userId}, \texttt{eventId}, \texttt{eventTitle}, čas
odoslania \texttt{scheduledAt} a príznak \texttt{sent}. Tento záznam sa
pri opustení udalosti alebo vymazaní eventu tiež odstráni. Perzistencia vo
Firestore slúži ako záloha pre prípad, keby sa lokálny alarm stratil
(napríklad po reštarte zariadenia).

